\documentclass[12pt]{article}
\usepackage{amsmath,amssymb}

\begin{document}
\section*{{2025 AMC 12A Problems/Problem 24}}
Source: AoPS Wiki

\subsection*{Solution}
Let the center of the large circle be $O$ and the centers of the $12$ circles be $A_1, A_2, A_3, \dots, A_{12}$ . Triangle $OA_1A_2$ has side lengths $r+1, r+1, 2$ , with the angle opposite $2$ being $360/12 = 30$ . Note that the line connecting $A_1$ and $A_2$ go through their common point of tangency, by definition, which causes $A_1A_2$ to have a length of $2$ . Drawing the angle bisector of the $30$ degree angle, we split $OA_1A_2$ into two congruent right triangles, each with hypotenuse $r+1$ and side opposite the $15$ degree angle $1$ . From here, note that $\sin{15} = \frac{\sqrt{6}-\sqrt{2}}{4}$ , which be derived using the trigonometric identity $\sin{(A-B)} = \sin{A} \cos{B} - \sin{B} \cos{A}$ , with $A=45$ and $B=30$ . In our right triangle, $\sin{15} = \frac{1}{r+1} = \frac{\sqrt{6}-\sqrt{2}}{4}$ . Let $x=r+1$ . Solving for $x$ , we get $x = \frac{4}{\sqrt{6}-\sqrt{2}}$ . Rationalizing, we get that $x = \sqrt{6}+\sqrt{2}$ . Remember $x = r+1 = \sqrt{6}+\sqrt{2}$ , so $r = \sqrt{6}+\sqrt{2} - 1$ . Therefore, our answer is $6+2-1 = \boxed{7}.$ ~lprado

\end{document}
